\section*{Chapitre \ref{ch:qm:convex-optimisation} --- Optimisation convexe}

\begin{solution}{exo:qm:convex-optimisation:1}
Oui~: $\sqrt{w^\top\Sigma w} = \lVert L^\top w\rVert$ est une norme composée avec une application linéaire, donc convexe. Non~: l'ensemble des portefeuilles dont la variance atteint un niveau donné est le complémentaire d'un ellipsoïde~; le
segment reliant $w$ et $-w$ passe par $0$, où la variance est nulle.
\end{solution}

\begin{solution}{exo:qm:convex-optimisation:2}
Au premier ordre~: $0.0159 \times 0.02 = 0.032\%$, soit environ $3.180\% + 0.032\% = 3.212\%$. Une nouvelle résolution à 22~\% donne 3.211~\%~: le multiplicateur surestime légèrement le gain, car le prix baisse à mesure que la
limite se relâche.
\end{solution}

\begin{solution}{exo:qm:convex-optimisation:3}
On introduit $u, v \ge 0$ avec $w - w_0 = u - v$ et on impose $\mathbf 1^\top(u + v) \le \tau$. À l'optimum, $u_iv_i = 0$ pour tout $i$ dès que la limite est saturée, donc $u + v = |w - w_0|$.
\end{solution}

\begin{solution}{exo:qm:convex-optimisation:4}
La stationnarité $\Sigma w + \nu\mathbf 1 = 0$ donne $w = -\nu\Sigma^{-1}\mathbf 1$~; la contrainte de budget donne $\nu = -1/(\mathbf 1^\top\Sigma^{-1}\mathbf 1)$, d'où $w = \Sigma^{-1}\mathbf 1/(\mathbf 1^\top\Sigma^{-1}\mathbf 1)$.
La valeur absolue du multiplicateur est la variance minimale, $1/(\mathbf 1^\top\Sigma^{-1}\mathbf 1)$~: c'est le taux auquel la demi-variance croît avec le budget.
\end{solution}

\begin{solution}{exo:qm:convex-optimisation:5}
$\lambda = (1, 1)$ sur les deux bornes supérieures et $\nu = -1$ sur la contrainte d'égalité~: $G^\top\lambda + A^\top\nu = (1, 1) - (1, 1) = 0$ et $h^\top\lambda + b^\top\nu = 0.5 - 1 = -0.5 < 0$. En clair~: $x_1 + x_2 \le 0.5$
contredit $x_1 + x_2 = 1$.
\end{solution}

\begin{solution}{exo:qm:convex-optimisation:6}
KKT~: $\alpha = 2\lambda\Sigma w$, donc $w \propto \Sigma^{-1}\alpha$~; la limite saturée fixe l'échelle, $w = \sigma\Sigma^{-1}\alpha/\sqrt{\alpha^\top\Sigma^{-1}\alpha}$, et la valeur optimale vaut $\sigma\sqrt{\alpha^\top\Sigma^{-1}\alpha}$,
dont la dérivée en $\sigma$ vaut $\sqrt{\alpha^\top\Sigma^{-1}\alpha}$, soit le \omterm{def:qm:estimation:sharpe}{ratio de Sharpe} maximal~: c'est le \omterm{def:qm:convex-optimisation:strong}{prix fictif} de la limite de volatilité.
\end{solution}

\begin{solution}{exo:qm:convex-optimisation:7}
\texttt{constraint\_costs()}~: le turnover coûte 0.44~\% de l'utilité, l'exposition brute 0.28~\%, les limites de position 0.19~\%, la neutralité sectorielle 0.014~\%, la neutralité en dollars 0. La contrainte en dollars est la somme des cinq contraintes
sectorielles, donc la retirer ne change rien~; son multiplicateur dans la solution renvoyée par le solveur n'est pas unique et ne dit rien.
\end{solution}

\begin{solution}{exo:qm:convex-optimisation:8}
Retirer des contraintes dans un ordre arbitraire revient à négocier un portefeuille qui viole une limite fixée par l'ordre du code, et non par le jugement de quiconque, et masque le conflit. Calculer le \omterm{def:qm:convex-optimisation:certificate}{certificat
d'infaisabilité}, rapporter les règles en conflit à leurs responsables, et ne passer aucun nouvel ordre (ou un repli documenté) tant qu'une limite n'a pas été modifiée délibérément.
\end{solution}

\begin{solution}{pb:qm:convex-optimisation:1}
\textbf{1.} Rendement espéré 3.78~\%, volatilité 3.47~\%, utilité 3.18~\% (annuels).
\textbf{2.} Exposition brute (100~\%), turnover (20~\%), 43 limites de position~; les neutralités sont vérifiées sous forme d'équations.
\textbf{3.} 18 itérations~; résidus de stationnarité et de complémentarité inférieurs à $10^{-9}$.
\textbf{4.} 0.0159~: 1~\% de turnover supplémentaire vaut 1.6 points de base d'utilité.
\textbf{5.} 0.0185 par unité d'exposition brute.
\textbf{6.} De 2.88~\% à 5~\%, à 3.18~\% à 20~\% et à 3.62~\% à 80~\%.
\textbf{7.} Ce sont les pentes~: 0.0210, 0.0184, 0.0159, 0.0119 et 0.0011 à 10~\%, 15~\%, 20~\%, 30~\% et 80~\%, qui coïncident avec les différences finies de la courbe.
\textbf{8.} Turnover 0.44~\%, exposition brute 0.28~\%, positions 0.19~\%, neutralité sectorielle 0.014~\%, neutralité en dollars 0.
\textbf{9.} La neutralité en dollars découle des neutralités sectorielles (de leur somme)~; elle est donc redondante~; les multiplicateurs des contraintes redondantes ne sont pas uniques.
\textbf{10.} Le rendement espéré passe de 3.60~\% à 3.92~\% entre 10~\% et 30~\% de turnover~: 1.6 points de base par 1~\% de turnover, proche du prix de l'utilité.
\textbf{11.} Les limites sont infaisables~: la phase 1 se termine à $t^* = 0.005$.
\textbf{12.} Pondérations $1/40$ sur chaque ligne de position minimale du secteur 0 et multiplicateur sur l'équation de neutralité du secteur~; la somme des lignes donne «~secteur 0 long d'au moins 20~\%~» face à «~secteur 0 exactement
0~»~: la combinaison $G^\top\lambda + A^\top\nu = 0$ avec $h^\top\lambda + b^\top\nu = -0.005 < 0$.
\textbf{13.} Sous forme de contrainte conique $\lVert L^\top w\rVert \le 0.04$~; son \omterm{def:qm:convex-optimisation:strong}{prix fictif} vaut 1.033 ($= \gamma \times 0.04$ avec $\gamma = 25.8$), soit le rendement espéré par unité de volatilité autorisée~; le portefeuille rapporte 5.31~\%.
\textbf{14.} Une matrice à valeurs propres négatives rend négative la variance de certains portefeuilles, que l'optimiseur trouvera et exploitera.
\textbf{15.} En 74 projections, elle supprime les 21 valeurs propres négatives à une distance de Frobenius de 4.45, sans modifier aucun élément de plus de 0.31.
\textbf{16.} Turnover~: c'est la contrainte la plus coûteuse (0.44~\% de l'utilité), et son prix est de 1.6 points de base par 1~\% échangé.
\textbf{17.} Elle ne se justifie que face aux coûts de transaction que la limite est censée épargner~: si échanger 1~\% du capital coûte moins de 1.6 points de base, la limite est trop serrée~; elle n'est qu'un substitut d'un modèle de coûts qui devrait
figurer dans la fonction objectif.
\textbf{18.} Statut, résidus KKT, valeur de chaque contrainte saturée et son multiplicateur, classement par relâchement, et tout certificat.
\textbf{19.} Résultat nommé~: \emph{la contrainte la plus coûteuse}~: le \omterm{def:qm:convex-optimisation:strong}{prix fictif} de la limite de turnover vaut 0.0159, soit 1.6 points de base d'utilité par 1~\% de turnover, et les contraintes se classent ainsi~: turnover (0.44~\%), exposition
brute (0.28~\%), limites de position (0.19~\%), neutralité sectorielle (0.014~\%), neutralité en dollars (redondante, 0).
\textbf{20.} Il chiffre chaque limite dans les unités de la fonction objectif, ce qui permet de débattre des limites avec des nombres, et il établit l'infaisabilité lorsqu'il n'y a rien à échanger.
\end{solution}

\begin{solution}{iq:qm:convex-optimisation:1}
Parce que tout optimum local est global~: un solveur qui satisfait les conditions d'optimalité locale tient la réponse, et ces conditions (KKT) se vérifient~; la dualité fournit alors des bornes, des prix
et des \omterm{def:qm:convex-optimisation:certificate}{certificats d'infaisabilité}, et il existe des algorithmes en temps polynomial.
\iqlookfor{Optimum local égal à optimum global, optimalité vérifiable et dualité.}
\end{solution}

\begin{solution}{iq:qm:convex-optimisation:2}
La variable duale optimale d'une contrainte~: le taux auquel l'objectif optimal s'améliore par unité de relâchement. Montrer au comité le multiplicateur de la limite de risque ou de turnover en rendement par
unité, le comparer à ce contre quoi la limite protège (coût, risque), et confirmer en résolvant à nouveau avec la limite relâchée.
\iqlookfor{Valeur marginale dans les unités de l'objectif, et confirmation par relâchement.}
\end{solution}

\begin{solution}{iq:qm:convex-optimisation:3}
Faisabilité du primal et du dual ($\lambda \ge 0$), stationnarité du \omterm{def:qm:convex-optimisation:dual}{lagrangien}, et $\lambda_ig_i(x) = 0$. Dans un portefeuille~: une limite qui n'est pas saturée a un prix nul, et une limite dont le prix est positif
est exactement sur sa borne~; les positions strictement à l'intérieur de leurs limites ont une utilité marginale égale aux prix des contraintes saturées qu'elles touchent.
\iqlookfor{Les quatre conditions et la lecture économique de la complémentarité.}
\end{solution}

\begin{solution}{iq:qm:convex-optimisation:4}
Résoudre un problème de phase 1 (ou auto-dual) et lire le certificat de Farkas~: les contraintes à poids positif forment un sous-ensemble en conflit, qui peut être signalé à leurs responsables. En variante,
minimiser la violation totale avec des poids reflétant les priorités.
\iqlookfor{Des certificats plutôt que des essais successifs.}
\end{solution}

\begin{solution}{iq:qm:convex-optimisation:5}
Les \omterm{def:qm:convex-optimisation:ipm}{méthodes de points intérieurs} donnent une grande précision en quelques dizaines d'itérations quelle que soit la taille, avec une structure exploitable (covariance factorielle) dans l'algèbre linéaire~; les méthodes du premier ordre, comme ADMM, ont
des itérations peu coûteuses, démarrent bien à chaud à partir de la solution de la veille et atteignent rapidement une précision suffisante pour le trading. Pour un rééquilibrage quotidien avec un \omterm{def:qm:covariance-estimation-and-random-matrices:factor}{modèle factoriel}, l'une ou l'autre fonctionne~; ADMM avec démarrage à chaud est attrayant
en cours de séance, les \omterm{def:qm:convex-optimisation:ipm}{méthodes de points intérieurs} lorsque les duaux doivent être précis.
\iqlookfor{Précision face aux démarrages à chaud, et exploitation de la structure factorielle.}
\end{solution}

\begin{solution}{iq:qm:convex-optimisation:6}
Tracking error~: $\lVert L^\top(w - b)\rVert_2 \le \mathrm{TE}$ avec $\Sigma = LL^\top$, un cône du second ordre. Impact $|x|^{3/2} \le s$~: on décompose $x$ en parties positives et on impose $x^2 \le rs$ et $r^2 \le x$ avec
$r \ge 0$, deux cônes tournés~; on somme les $s$ dans la fonction objectif.
\iqlookfor{L'astuce de Cholesky et la décomposition en cône de puissance.}
\end{solution}
